
\documentclass{article}
\usepackage{colortbl}
\usepackage{makecell}
\usepackage{multirow}
\usepackage{supertabular}

\begin{document}

\newcounter{utterance}

\centering \large Interaction Transcript for game `cladder', experiment `full\_v1.5\_default', episode 2679 with score{-}Qwen3.5{-}9B{-}sp.
\vspace{24pt}

{ \footnotesize  \setcounter{utterance}{1}
\setlength{\tabcolsep}{0pt}
\begin{supertabular}{c@{$\;$}|p{.15\linewidth}@{}p{.15\linewidth}p{.15\linewidth}p{.15\linewidth}p{.15\linewidth}p{.15\linewidth}}
    \# & \multicolumn{2}{c}{Player} && \multicolumn{2}{c}{Game Master} \\
    \hline

    \theutterance \stepcounter{utterance}  
    & & & \multicolumn{4}{p{0.6\linewidth}}{
        \cellcolor[rgb]{0.9,0.9,0.9}{
            \makecell[{{p{\linewidth}}}]{
                \texttt{\tiny{[P1$\langle$GM]}}
                \texttt{You are an expert in causal inference. The following question is not a typical commonsense query, but rather a meticulously designed question created by a professor specializing in causal inference, intended to assess the students' mastery of the course content.} \\
\texttt{Imagine a self{-}contained, hypothetical world with only the following conditions, and without any unmentioned factors or causal relationships: Gene has a direct effect on smoking and lung cancer. Pollution has a direct effect on smoking and lung cancer. Smoking has a direct effect on lung cancer. Pollution is unobserved. For people with nonsmoking genes and nonsmokers, the probability of lung cancer is 56\%. For people with nonsmoking genes and smokers, the probability of lung cancer is 42\%. For people with smoking genes and nonsmokers, the probability of lung cancer is 60\%. For people with smoking genes and smokers, the probability of lung cancer is 40\%. For people with nonsmoking genes and with low pollution, the probability of smoking is 76\%. For people with nonsmoking genes and with high pollution, the probability of smoking is 48\%. For people with smoking genes and with low pollution, the probability of smoking is 46\%. For people with smoking genes and with high pollution, the probability of smoking is 21\%. The overall probability of high pollution is 41\%. Does gene positively affect lung cancer through smoking?} \\
\texttt{Start your answer with 'Yes' or 'No', followed by additional reasoning or evidence to support your explanation.} \\
            }
        }
    }
     \\ \\

    \theutterance \stepcounter{utterance}  
    & \multicolumn{4}{p{0.6\linewidth}}{
        \cellcolor[rgb]{0.561,0.737,0.561}{
            \makecell[{{p{\linewidth}}}]{
                \texttt{\tiny{[P1$\rangle$GM]}}
                \texttt{No.} \\
\\ 
\texttt{**Reasoning:**} \\
\\ 
\texttt{To determine if the gene positively affects lung cancer *through* smoking, we must evaluate the **causal effect of smoking on lung cancer conditional on the gene**. In a causal graph where \$G \textbackslash{}to S\$, \$G \textbackslash{}to L\$, \$P \textbackslash{}to S\$, \$P \textbackslash{}to L\$, and \$S \textbackslash{}to L\$, the path from Gene to Lung Cancer via Smoking is \$G \textbackslash{}to S \textbackslash{}to L\$. For this path to represent a positive effect, individuals with a specific gene who smoke must have a higher probability of lung cancer than those with the same gene who do not smoke, holding the gene constant.} \\
\\ 
\texttt{Let us calculate the conditional probability of lung cancer given the gene and smoking status using the provided data:} \\
\\ 
\texttt{1.  **For individuals with Nonsmoking Genes (\$G=0\$):**} \\
\texttt{    *   \$P(L=1 | G=0, S=0) = 56\textbackslash{}\%\$} \\
\texttt{    *   \$P(L=1 | G=0, S=1) = 42\textbackslash{}\%\$} \\
\texttt{    *   Effect of smoking: \$42\textbackslash{}\% {-} 56\textbackslash{}\% = {-}14\textbackslash{}\%\$} \\
\\ 
\texttt{2.  **For individuals with Smoking Genes (\$G=1\$):**} \\
\texttt{    *   \$P(L=1 | G=1, S=0) = 60\textbackslash{}\%\$} \\
\texttt{    *   \$P} \\
            }
        }
    }
    & & \\ \\

    \theutterance \stepcounter{utterance}  
    & & & \multicolumn{2}{p{0.3\linewidth}}{
        \cellcolor[rgb]{0.9,0.9,0.9}{
            \makecell[{{p{\linewidth}}}]{
                \texttt{\tiny{[GM$|$GM]}}
                \texttt{No} \\
            }
        }
    }
    & & \\ \\

    \theutterance \stepcounter{utterance}  
    & & & \multicolumn{2}{p{0.3\linewidth}}{
        \cellcolor[rgb]{0.9,0.9,0.9}{
            \makecell[{{p{\linewidth}}}]{
                \texttt{\tiny{[GM$|$GM]}}
                \texttt{yes} \\
            }
        }
    }
    & & \\ \\

    \theutterance \stepcounter{utterance}  
    & & & \multicolumn{2}{p{0.3\linewidth}}{
        \cellcolor[rgb]{0.9,0.9,0.9}{
            \makecell[{{p{\linewidth}}}]{
                \texttt{\tiny{[GM$|$GM]}}
                \texttt{game\_result = LOSE} \\
            }
        }
    }
    & & \\ \\

\end{supertabular}
}

\end{document}
